\documentclass[../../../../main.tex]{subfiles}
\begin{document}

平面 $S$ の一点 $A$ と正数 $\alpha$ $(\alpha<180)$ をとる．
点の集合としての $S$ から $S$ への写像 $\varphi$ が，
次の三つの条件 (i)，(ii)，(iii) をみたすとき，
$\varphi$ は $A$ を中心とする正の向きの $\alpha^\circ$ 回転と呼ばれる．

\begin{enumerate}
  \item $\varphi(A)=A$，
  \item $S$ の任意の点 $P$ $(\neq A)$ に対し，$AP=A\varphi(P)$，$\angle PA\varphi(P)=\alpha^\circ$，
  \item 人が三角形 $P\varphi(P)A$ の周を一周し，$P$，$\varphi(P)$，$A$ の順序に頂点を通るとき，
  三角形の内部は常に人の左側にある．
\end{enumerate}

いま $S$ 上に相異なる二点 $A$，$B$ をとり，
$A$ を中心とする正の向きの $60^\circ$ 回転を $f$，
$B$ を中心とする正の向きの $60^\circ$ 回転を $g$ とする．
これに対し，$f$ と $g$ の合成写像 $h=g \circ f$ が，$h(P)=g(f(P))$ によって定義される．

\begin{figure}[htbp]
  \centering
  \begin{tikzpicture}[scale=1.5, >=stealth]
    \coordinate (A) at (0,0);
    \coordinate (P) at (25:2);
    \coordinate (FP) at (85:2.2);

    \draw[->, thick] (A) -- (P) node[right] {$P$};
    \draw[->, thick] (A) -- (FP) node[above] {$f(P)$};

    \draw[->] (A) +(25:0.6) arc (25:85:0.6);
    \node at ($(A)+(55:0.85)$) {$60^\circ$};

    \node[below left] at (A) {$A$};
  \end{tikzpicture}
\end{figure}

\begin{enumerate}
  \item このとき，点 $h(A)$ と $h(B)$ は，$A$，$B$ に対して，どのような位置にあるかを求め，図示せよ．
  \item $h$ はある点 $O$ を中心とする正の向きの回転であることを示し，点 $O$ および回転角を求めよ．
\end{enumerate}

\end{document}
